\documentclass[10pt,a4paper]{amsart}
\usepackage[margin=19mm]{geometry}
\usepackage{amsmath,amssymb,mathtools}
\usepackage[scr=rsfs,cal=euler]{mathalfa}
\usepackage{xfrac}
\usepackage[unicode,hidelinks,bookmarks=false]{hyperref}
\newcommand{\ie}{i.e.\ }
\newcommand{\cf}{cf.\ }
\newcommand{\resp}{respectively\ }
\newcommand{\Subsection}{\subsection}
\providecommand{\nobreakdash}{\nobreak\hbox{-}}
\setlength{\emergencystretch}{2em}
\multlinegap0pt
\usepackage{amssymb}
\usepackage[shortlabels]{enumitem}
\newtheorem{prop}{Proposition}
\def\Rm{\operatorname{Rm}}
\def\tr{\operatorname{tr}}
\title{Asymptotic Profiles and Non-Trivial Breathers in Kähler-Ricci Flow}
\author{Longteng Chen}
\date{Source: arXiv:2606.08837v1 (CC BY 4.0)}
\begin{document}
\maketitle

\noindent\emph{Source and licence.} Asymptotic Profiles and Non-Trivial Breathers in Kähler-Ricci Flow. Version arXiv:2606.08837v1 (CC BY 4.0), distributed by arXiv under the Creative Commons Attribution 4.0 International licence, http://creativecommons.org/licenses/by/4.0/, as recorded in the arXiv licence metadata for this exact version and shown on the abstract page https://arxiv.org/abs/2606.08837v1. Full licence notice: \texttt{licenses/arxiv-2606.08837-CC-BY-4.0.txt}.\par\smallskip

\noindent\textit{Editorial scope.} Counted unit: the article's prop estimates (source0.tex lines 661--668). The article's complete original proof of it is delivered with it (source0.tex lines 669--739, one \texttt{proof} environment of 71 lines). No mathematics is written, repaired or re-derived: the statement and the proof are the article's own lines, in the article's own notation, and no retained line is substituted or re-flowed.  Retained uncounted as the source-local context the proof needs: lines 551--563.  Nothing was omitted from any retained span, so the package carries no omission record.  No retained line contains a citation, so the wrapper carries no bibliography and no bibliography entry.  No retained line contains a figure environment, an image-inclusion command or drawing code, so no image is delivered and the package declares no asset dependency.  Editorial formatting only, in addition: an AMS \texttt{amsart} preamble that restates the article's own declarations and macros used by the retained lines, namely the environments \texttt{prop} on separate counters, together with the standard packages those lines need.  The retained environments print in this document as Proposition 1 in order of appearance, whereas the article prints the same items with the numbering of its own full document.  The complete native source of the article as distributed in the arXiv e-print for this exact version is bundled unchanged as \texttt{sources/202360/source0.tex}.
\begin{prop}[Quantification of initial data]\label{quantification of initial data}
    There exists a real-valued $JX$-invariant function $\psi_0\in C^\infty(M)$ such that the following holds:
    \begin{enumerate}
        \item The function $\psi_0$ is the K\"ahler potential of $\omega_0$, i.e, $\omega_0=\omega+i\partial\bar\partial\psi_0$.
        \item Let $f$ be the normalized soliton potential, there exists a constant $C>0$ such that on $M$,
        \begin{equation*}
            \begin{split}
               & (e^{-a}-1)f-C\le\psi_0\le (e^a-1)f+C;\\
               & (e^{-a}-1)f-C\le\frac{X}{2}\cdot\psi_0\le (e^a-1)f+C.
            \end{split}
        \end{equation*}
    \end{enumerate}
\end{prop}

\begin{prop}\label{rough estimates}
    We can take $\lambda>>1$ sufficiently large such that on $\{(x,t)\ |\ r(x)^2\ge\lambda(t+1),\ t\in[0,T)\}$, we have that
    \begin{enumerate}
        \item $e^{-a}g(t+1)\le g_\varphi(t)\le e^ag(t+1)$;
        \item $(e^{-a}-1)f-C(t+1)\le\varphi(t)\le (e^a-1)f+C(t+1)$;
        \item $(e^{-a}-1)f-C(t+1)\le\frac{X}{2}\cdot\varphi(t)\le (e^a-1)f+C(t+1)$.
    \end{enumerate}
\end{prop}

\begin{proof}
    First take $\lambda>>1$ such that on $\{r^2\ge\lambda\}$, there is an $\varepsilon>0$ such that
    \begin{equation*}
       e^{-a+\varepsilon}g \le g_0\le e^{a-\varepsilon}g.   
    \end{equation*}
Since on $\{r^2\ge\lambda(t+1)\}$, we have $|\Rm(g_\varphi(t))|_{g_\varphi(t)}\le \frac{C}{r^2}$, it follows that
\begin{equation*}
    e^{-\frac{Ct}{r^2}}g_0\le g_\varphi(t)\le e^{\frac{Ct}{r^2}}g_0.
\end{equation*}
Similarly, we have
\begin{equation*}
    e^{-\frac{Ct}{r^2}}g\le g(t+1)\le e^{\frac{Ct}{r^2}}g.
\end{equation*}
Combine these, we get
\begin{equation*}
   \begin{split}
        &g_\varphi(t)\le e^{a-\varepsilon+\frac{2Ct}{r^2}}g\le e^{a-\varepsilon+\frac{2C}{\lambda}}g;\\
        &g_\varphi(t)\ge e^{-a+\varepsilon-\frac{2Ct}{r^2}}g(t+1)\le e^{-a+\varepsilon-\frac{2C}{\lambda}}g.
   \end{split}
\end{equation*}
Take $\lambda>>1$ such that $\varepsilon-\frac{2C}{\lambda}>0$, then we have 
\begin{equation*}
    e^{-a}g(t+1)\le g_\varphi(t)\le e^ag(t+1),
\end{equation*}
holds on $\{r^2\ge\lambda(t+1)\}$.
   
    Recall the equation of Monge-Amp\`ere flow:
    \begin{equation*}
        \frac{\partial}{\partial t}\varphi=\log\frac{\omega_\varphi(t)^n}{\omega(t+1)^n}.
    \end{equation*}
    By integrating, we get
    \begin{equation*}
        |\varphi(t)-\varphi(0)|\le na t.
    \end{equation*}
    Recall that $\varphi(0)=\psi_0\le (e^a-1)f+C$ as in Proposition \ref{quantification of initial data}, then we have
    \begin{equation*}
         \varphi(t)\le (e^a-1)f+C+nat.
    \end{equation*}
Similarly, it follows that
\begin{equation*}
    \varphi(t)\ge (e^{-a}-1)f-C-nat.
\end{equation*}
    We also have
    \begin{equation*}
        \begin{split}
            \frac{X}{2}\cdot\frac{\partial}{\partial t}\varphi&=\tr_{\omega_\varphi(t)}\mathcal{L}_{\frac{X}{2}}\omega_\varphi(t)-\tr_{\omega(t)}\mathcal{L}_{\frac{X}{2}}\omega(t+1)\\
            &=\tr_{\omega_\varphi(t)}\mathcal{L}_{\frac{X}{2}}\omega(t+1)-\tr_{\omega(t)}\mathcal{L}_{\frac{X}{2}}\omega(t+1)+\tr_{\omega_\varphi(t)}\mathcal{L}_{\frac{X}{2}}i\partial\bar\partial\varphi(t).
        \end{split}
    \end{equation*}
    Since we can control $\tr_{\omega_\varphi(t)}\mathcal{L}_{\frac{X}{2}}i\partial\bar\partial\varphi(t)$ in the following way: \begin{equation*}
    \begin{split}
            |\tr_{\omega_\varphi(t)}\mathcal{L}_{\frac{X}{2}}i\partial\bar\partial\varphi(t)|&\le C(n)|g_\varphi(t)|_{g(t+1)}\left(|X|_{g(t)}|\nabla^{g(t+1)}i\partial\bar\partial\varphi(t)|_{g(t+1)}+|\nabla^{g(t+1)}X|_{g(t+1)}|i\partial\bar\partial\varphi(t)|_{g(t+1)}\right)\\
            &\le C.
    \end{split}
    \end{equation*}
    Therefore, $ |\frac{X}{2}\cdot\frac{\partial}{\partial t}\varphi|$ is uniformly bounded by some constant $C_1>0$ on $\{r^2\ge\lambda(t+1)\}$. By integration, we have
    \begin{equation*}
         |\frac{X}{2}\cdot\varphi(t)-\frac{X}{2}\cdot\varphi(0)|\le C_1t.
    \end{equation*}
    Recall that $\frac{X}{2}\cdot\varphi(0)=\frac{X}{2}\cdot\psi_0\le (e^a-1)f+C$. Then we have
    \begin{equation*}
        \frac{X}{2}\cdot\varphi(t)\le (e^a-1)f+C+C_1t;\quad \frac{X}{2}\cdot\varphi(t)\ge (e^{-a}-1)f-C-C_1t.
    \end{equation*}
    Take $C=C+C_1+na$, we have
    \begin{equation*}
        \begin{split}
           & (e^{-a}-1)f-C(t+1)\le \varphi(t)\le (e^a-1)f+C(t+1);\\
           &(e^{-a}-1)f-C(t+1)\le\frac{X}{2}\cdot\varphi(t)\le (e^a-1)f+C(t+1).  
        \end{split}
          \end{equation*}
\end{proof}

\end{document}
